\section{Contract-Aware Update Analysis}
\label{sec:gate}

This section develops the update conditions used by the empirical studies.
The current- and future-history gates distinguish action filtering from
certified learning.  The reward-influence analysis then characterizes how a
bounded reward edit can change a linear release-gate margin within one policy-
gradient batch.  The evaluated multi-batch attack remains the fixed
target-shaping rule described in Section~\ref{sec:study}.

\subsection{Current- and Future-History Gates}

The certificate condition has two levels.  The current-history condition
requires the updated action at the observed history to remain certified.  The
lifecycle condition extends this requirement to future attacked histories
reachable after the update.

\begin{lemma}[Current-history gate]
\label{lem:current-gate}
Let $M=(\pi_\theta,U,\mathsf{Cert})$ be a history-dependent online controller.
For an attacked history $h$, let $X_{\mathcal A}(h)$ be the plausible-state set
and $K_\Gamma(h)$ its certified kernel.  If $M$ accepts a sound certificate for
$\pi_{\theta_{t+1}}$, where $\theta_{t+1}=U(\theta_t,h)$, then
\begin{equation}
    \pi_{\theta_{t+1}}(h)\in K_\Gamma(h).
\end{equation}
Consequently, if $K_\Gamma(h)=\emptyset$, then $M$ cannot both accept the certificate and allow a nontrivial online update at $h$.
\end{lemma}

\begin{proof}
The certified action must satisfy the certificate for every state compatible
with the attacked history.  These states form $X_{\mathcal A}(h)$, and their
common certified actions form $K_\Gamma(h)$.  Continued certified operation
therefore requires $\pi_{\theta_{t+1}}(h)\in K_\Gamma(h)$.  An empty kernel
requires rejection, freeze, rollback, or trusted information that changes the
plausible-state set.
\end{proof}

\begin{proposition}[Action filtering is not update certification]
\label{prop:action-filter-insufficient}
Let $F(v,h)$ be an action filter that returns an applied input in
$K_\Gamma(h)$ whenever the kernel is nonempty.  Consider a mechanism that
first computes $\theta^+=U(\theta_t,h)$ and then applies
\begin{equation}
    u_t=F(\pi_{\theta^+}(h),h),
\end{equation}
The mechanism accepts the update without constraining the raw updated policy
$\pi_{\theta^+}$.  If $\pi_{\theta^+}(h)\notin K_\Gamma(h)$, the applied input
can be certified while the updated learner remains uncertified.  Action-only
filtering therefore does not imply Lemma~\ref{lem:current-gate}.
\end{proposition}

\begin{proof}
The filter property gives $u_t\in K_\Gamma(h)$, so the current actuator input
can satisfy the certificate.  Lemma~\ref{lem:current-gate} instead constrains
the updated controller through
$\pi_{\theta^+}(h)\in K_\Gamma(h)$.  When this membership fails, the
certificate covers only the filtered actuator channel.  Accepting the update
under that certificate creates a stale certified policy.
\end{proof}

The lemma and proposition define local gates.  A policy update also determines
actions at future histories.  Let $\mathcal H_{\mathcal A}(h,\theta^+)$ denote
the attacked histories reachable after $h$ under $\pi_{\theta^+}$ and attack
family $\mathcal A$.  Define the certified parameter set
\begin{equation}
\Theta_\Gamma(h)=
\{\theta:\forall h'\in\mathcal H_{\mathcal A}(h,\theta),
\pi_\theta(h')\in K_\Gamma(h')\}.
\end{equation}

\begin{theorem}[Future-history lifecycle gate]
\label{thm:gate}
Let $\theta^+=U(\theta_t,h)$ be a candidate online update.  Suppose $M$
accepts a sound lifecycle certificate for $\pi_{\theta^+}$ from history $h$
over attack family $\mathcal A$.  Then
\begin{equation}
    \theta^+\in\Theta_\Gamma(h),
\end{equation}
equivalently,
\begin{equation}
    \forall h'\in\mathcal H_{\mathcal A}(h,\theta^+),\quad
    \pi_{\theta^+}(h')\in K_\Gamma(h').
\end{equation}
If membership fails for a reachable $h'$, accepting the update creates a
stale certified policy.  If $K_\Gamma(h')=\emptyset$, certified learning must
reject, freeze, roll back, or obtain trusted information that changes
$X_{\mathcal A}(h')$.
\end{theorem}

\begin{proof}
A lifecycle certificate covers every attacked history reachable under the
updated controller and admissible attack family.  Applying
Lemma~\ref{lem:current-gate} pointwise gives
$\pi_{\theta^+}(h')\in K_\Gamma(h')$ for each such history.  This is the
definition of $\theta^+\in\Theta_\Gamma(h)$.  A failed membership leaves the
updated action map uncertified at that history.  An empty kernel additionally
requires a changed information set or rejection of certified operation.
\end{proof}

\begin{proposition}[Linear-policy parameter gate]
\label{prop:parameter-gate}
Consider a finite abstraction of future attacked histories indexed by
$z\in\mathcal Z$ and a linear policy
$\pi_\theta(z)=\theta^\top\phi(z)$.  Suppose each certificate produces an
interval kernel $K_\Gamma(z)=[\ell_z,u_z]$.  The lifecycle condition over
$\mathcal Z$ is then equivalent to
\begin{equation}
    \theta^\top\phi(z)\le u_z,\qquad
    -\theta^\top\phi(z)\le -\ell_z,\qquad
    \forall z\in\mathcal Z.
\end{equation}
Optional box constraints on $\theta$ preserve a convex polytope.  Rejecting
candidates outside this set, or projecting onto it, is sound for the finite
abstraction.
\end{proposition}

\begin{proof}
For fixed $z$, membership $\pi_\theta(z)\in[\ell_z,u_z]$ is equivalent to the
two inequalities above.  Enforcing them for every $z\in\mathcal Z$ implements
the lifecycle gate on the abstraction.  Parameter box constraints preserve
convexity.  A feasible or projected candidate satisfies every certified
interval kernel in the abstraction.
\end{proof}

\subsection{Normalized Reward Influence on a Release Gate}
\label{sec:reward-geometry}

We now connect the reward-record attacker to the halfspaces above.  Consider a
fixed on-policy batch of length $T$.  Let $J\in\mathbb R^{T\times d}$ contain
the policy-score row vectors, let $r\in\mathbb R^T$ be the logged rewards, and
let $H_{ij}=\mathbf 1[j\ge i]$ be the reward-to-go operator.  Define the
centering matrix
\begin{equation}
C=I-\frac{1}{T}\mathbf 1\mathbf 1^\top
\end{equation}
and let the attacker choose
$\delta\in[-\epsilon,\epsilon]^T$.  The centered return vector is
\begin{equation}
y(\delta)=CH(r+\delta).
\label{eq:centered-return}
\end{equation}
The fixed-batch assumption is appropriate for a reward-record-only attacker:
changing a reward after collection does not change the states, actions, or
score rows already in the batch.

\begin{proposition}[Standardized reward-influence set]
\label{prop:normalized-reward-influence}
Suppose the learner standardizes the centered reward-to-go advantages with
population standard deviation whenever $\|y(\delta)\|_2>0$.  Before optimizer
clipping, a gradient-ascent proposal of size $\alpha$ is
\begin{equation}
\widetilde\theta^+(\delta)=\theta+
\frac{\alpha}{\sqrt T}J^\top
\frac{y(\delta)}{\|y(\delta)\|_2}.
\label{eq:normalized-update}
\end{equation}
Consequently, the raw one-batch proposal set is the image of the centered
reward-to-go zonotope under normalization and the policy-score map; in general,
it is not a parameter-space zonotope.
\end{proposition}

\begin{proof}
The standardized advantage vector is
$y/\sqrt{T^{-1}\|y\|_2^2}=\sqrt T\,y/\|y\|_2$.  Substitution into the usual
mean score-function gradient
$T^{-1}J^\top(\sqrt T\,y/\|y\|_2)$ gives
\eqref{eq:normalized-update}.  Equation~\eqref{eq:centered-return} is affine
in the reward-box variable, but division by its Euclidean norm is not;
hence the normalized image is generally non-affine.
\end{proof}

\begin{theorem}[Positive support under global gradient clipping]
\label{thm:reward-halfspace-support}
Let $B=J^\top/\sqrt T$, let the learner cap the global gradient norm at
$G>0$, and assume coordinatewise parameter clipping is inactive over the
reward box.  The post-cap update before parameter clipping is
\begin{equation}
 \theta_G^+(y)=\theta+
 \frac{\alpha By}{\max\{\|y\|_2,\|By\|_2/G\}}.
\end{equation}
For a release halfspace $a^\top P\theta\le b$, define
$q=\alpha B^\top P^\top a$, $y_0=CHr$, and $M=CH$.  For any fixed $t>0$, a
reward perturbation whose gate-coordinate increment is at least $t$ exists if
and only if the following second-order-cone problem is feasible:
\begin{equation}
\begin{split}
 \lambda&\ge0,\qquad
 -\epsilon\lambda\mathbf1\le v\le\epsilon\lambda\mathbf1,\\
 \bar y&=\lambda y_0+Mv,\qquad q^\top\bar y=1,\\
 t\|\bar y\|_2&\le1,\qquad
 (t/G)\|B\bar y\|_2\le1.
\end{split}
\label{eq:soc-support-feasibility}
\end{equation}
Feasibility is monotone in $t$, so bisection on
$[0,\min\{\|q\|_2,\alpha G\|P^\top a\|_2\}]$ computes the exact positive
support
\begin{equation}
\max\left\{0,\sup_{\|\delta\|_\infty\le\epsilon}
\frac{q^\top y(\delta)}
{\max\{\|y(\delta)\|_2,\|By(\delta)\|_2/G\}}\right\}.
\end{equation}
The construction remains valid when the centered-return zonotope contains the
origin; if $G=\infty$, the second norm constraint is omitted.  If this support
does not exceed the current halfspace margin, no admissible reward edit can
cross that halfspace in the batch.
\end{theorem}

\begin{proof}
Global norm clipping gives the stated maximum in the denominator.  The ratio
is positively homogeneous, so optimization over the zonotope
$Z=\{y_0+M\delta:\|\delta\|_\infty\le\epsilon\}$ is equivalent to optimization
over its conic hull.  The first line of~\eqref{eq:soc-support-feasibility} is
the perspective representation of that hull: $v=\lambda\delta$.  This is the
standard homogenization used in fractional and conic
optimization~\cite{schaible_fractional_duality_1976,bauschke_homogenization_2023}.

If a positive ratio reaches $t$, rescale its conic representative until
$q^\top\bar y=1$; its two denominator branches are then exactly the last line
of~\eqref{eq:soc-support-feasibility}.  Conversely, any feasible $\bar y$ is
nonzero because of the unit-numerator equality and recovers a reward-box ray
whose ratio reaches $t$.  Thus the origin cannot create the degenerate
feasibility present in a direct zonotope formulation.  Each fixed-$t$ problem
is a second-order feasibility-cone problem~\cite{belloni_soc_feasibility_2008};
monotonicity permits quasiconvex bisection~\cite{agrawal_dqcp_2020}.
Cauchy--Schwarz supplies both stated upper bounds.  The margin result follows
by adding the maximum increment to the current gate coordinate.
\end{proof}

The theorem applies to one fixed batch.  Later batches change $J$, $r$, and
the visited states, producing a sequence of state-dependent normalized
updates.  The result covers global Euclidean gradient clipping.  Coordinatewise
parameter clipping remains outside the exact support computation.  The
implementation also skips standardization below a numerical tolerance, so
each reported cone witness must remain above that threshold.

\subsection{A Margin-Robust Finite-Abstraction Lift}

Proposition~\ref{prop:parameter-gate} is exact on its listed grid.  Lifting
that result to continuous reachable histories requires coverage and regularity
assumptions.  Let $z=\psi(h')$ be a finite-dimensional representation on
which the updated policy depends.  Suppose the fixed-controller certificate
can be written as finitely many obligations
\begin{equation}
 c_j(x,u,w,x^+)\le 0,\qquad j=1,\ldots,m,
\label{eq:certificate-obligation}
\end{equation}
where $x^+=f(x,u,w)$.  Invariant, discrete-time barrier, and bounded-horizon
reachability checks can be written in this form after augmenting the state
with time or certificate memory.

\begin{theorem}[Margin-robust finite-abstraction sufficiency]
\label{thm:finite-lift}
Fix a candidate snapshot $\pi_\theta$ and a future-history horizon.  Let
$\{(z_i,x_{iq},w_{ir})\}$ be a finite set such that, for every reachable
attacked history $h'$, every $x\in X_{\mathcal A}(h')$, and every
$w\in\mathcal W$, some sampled triple satisfies
\begin{equation}
\|z-z_i\|\le\delta_z,\quad
\|x-x_{iq}\|\le\delta_x,\quad
\|w-w_{ir}\|\le\delta_w,
\label{eq:joint-cover}
\end{equation}
where $z=\psi(h')$.  Assume:

\begin{enumerate}
    \item $\pi_\theta(z)\in\mathcal U$ and
    $\|\pi_\theta(z)-\pi_\theta(z')\|\le L_\pi\|z-z'\|$;
    \item the certificate model $\widehat f$ has uniform error
    $\|f(x,u,w)-\widehat f(x,u,w)\|\le\varepsilon_f$;
    \item $\widehat f$ is separately Lipschitz with constants
    $L_f^x,L_f^u,L_f^w$; and
    \item each $c_j$ is separately Lipschitz in
    $(x,u,w,x^+)$ with constants
    $L_j^x,L_j^u,L_j^w,L_j^+$.
\end{enumerate}

Define
\begin{equation}
\begin{split}
\eta_j={}&(L_j^x+L_j^+L_f^x)\delta_x\\
&+(L_j^u+L_j^+L_f^u)L_\pi\delta_z\\
&+(L_j^w+L_j^+L_f^w)\delta_w+L_j^+\varepsilon_f .
\end{split}
\label{eq:lift-margin}
\end{equation}
If every sampled triple satisfies
\begin{equation}
c_j\!\left(x_{iq},\pi_\theta(z_i),w_{ir},
\widehat f(x_{iq},\pi_\theta(z_i),w_{ir})\right)
\le-\eta_j
\label{eq:tightened-obligation}
\end{equation}
for all $j$, then $\pi_\theta(\psi(h'))\in K_\Gamma(h')$ for every reachable
attacked history $h'$ in the horizon.  Hence the snapshot satisfies the
continuous future-history lifecycle condition of
Theorem~\ref{thm:gate}.
\end{theorem}

\begin{proof}
Fix $h'$, $x\in X_{\mathcal A}(h')$, and $w\in\mathcal W$, and select a
sampled triple using~\eqref{eq:joint-cover}.  Write
$u=\pi_\theta(z)$ and $u_i=\pi_\theta(z_i)$.  Policy regularity gives
$\|u-u_i\|\le L_\pi\delta_z$.  The model-error and Lipschitz assumptions give
\begin{equation}
\begin{split}
\|f(x,u,w)-\widehat f(x_{iq},u_i,w_{ir})\|
\le{}&\varepsilon_f+L_f^x\delta_x\\
&+L_f^uL_\pi\delta_z+L_f^w\delta_w .
\end{split}
\end{equation}
Applying the four Lipschitz bounds of $c_j$ and collecting terms yields
\begin{equation}
\begin{split}
c_j(x,u,w,f(x,u,w))
\le{}&c_j(x_{iq},u_i,w_{ir},\\
&\widehat f(x_{iq},u_i,w_{ir}))+\eta_j\\
\le{}&0.
\end{split}
\end{equation}
The argument holds for every plausible state, disturbance, certificate
obligation, and reachable attacked history, which is precisely membership in
the corresponding kernels over the stated horizon.
\end{proof}

Theorem~\ref{thm:finite-lift} converts a justified joint cover into a
verification result.  Establishing that cover remains a separate obligation;
a finite rollout sample is insufficient.  The coverage audit in
Section~\ref{sec:study} provides empirical evidence about this premise without
constituting a cover proof.

\begin{proposition}[Plausible-set and freeze-frontier monotonicity]
\label{prop:plausible-monotonicity}
Fix a history domain $\mathcal H$, dynamics, certificate obligations, and
policy class.  If two information contracts satisfy
$X_1(h)\subseteq X_2(h)$ for every $h\in\mathcal H$, let
$K_{\Gamma,k}$ be the kernel induced by $X_k$ and let
$\Theta_{\Gamma,k}=\{\theta:\pi_\theta(h)\in K_{\Gamma,k}(h),
\ \forall h\in\mathcal H\}$.  Then
\begin{equation}
K_{\Gamma,2}(h)\subseteq K_{\Gamma,1}(h)
\quad\text{and}\quad
\Theta_{\Gamma,2}\subseteq\Theta_{\Gamma,1}.
\end{equation}
Consequently, for a nested ambiguity family
$X_{\rho_1}(h)\subseteq X_{\rho_2}(h)$ whenever $\rho_1\le\rho_2$, certified
update feasibility is non-increasing in $\rho$: once the certified parameter
set is empty, it remains empty at larger ambiguity.  Conversely, a trusted
anchor that pointwise shrinks the plausible-state sets can only weakly enlarge
the kernels and certified parameter set.
\end{proposition}

\begin{proof}
An action certified for every state in the larger set $X_2(h)$ is certified
for every state in its subset $X_1(h)$, proving kernel inclusion.  Requiring a
policy action to lie in the smaller kernel $K_{\Gamma,2}(h)$ for every
$h\in\mathcal H$ is at least as restrictive as requiring membership in
$K_{\Gamma,1}(h)$, proving parameter-set inclusion.  The frontier and anchor
statements follow by applying these inclusions to nested families.
\end{proof}

The monotonicity comparison holds on a fixed history domain and certificate
contract.  If an anchor changes which histories are reachable, or if two
sampled plausible sets are not nested, neither conclusion follows
automatically.  The P3 anchor audit therefore constructs nested samples
explicitly.

The three results have distinct scopes.  The future-history theorem states a
necessary certificate condition.  Proposition~\ref{prop:parameter-gate} is
exact on its finite abstraction.  Theorem~\ref{thm:finite-lift} supplies a
sufficient lift under its cover, regularity, and margin assumptions.

\subsection{Instantiating the Kernel}

For an invariant certificate $I\subseteq S$, the kernel is
\begin{equation}
\begin{aligned}
K_I(h)=\{u\in\mathcal U:\;&\forall x\in X_{\mathcal A}(h),
\forall w\in\mathcal W,\\
& f(x,u,w)\in I\}.
\end{aligned}
\end{equation}
For a barrier certificate $B(x)\ge 0$ with safe set
$C=\{x:B(x)\ge 0\}$, the corresponding kernel is
\begin{equation}
\begin{aligned}
K_B(h)=\{u\in\mathcal U:\;&\forall x\in X_{\mathcal A}(h),
\forall w\in\mathcal W,\\
& B(f(x,u,w))\ge 0\}.
\end{aligned}
\end{equation}
A reachability certificate starts from $X_{\mathcal A}(h)$ rather than a
nominal state estimate.  A runtime shield must select an action in
$K_\Gamma(h)$.  An empty kernel requires rejection, fail-safe operation, or a
trusted state anchor.

\subsection{A Minimal Counterexample}

Consider
\begin{equation}
    x_{t+1}=\lambda x_t+u_t,\qquad |u_t|\le u_{\max},\qquad S=[-1,1].
\end{equation}
Suppose FDI makes the same history compatible with $x_t=r$ and $x_t=-r$,
where $r\in(0,1]$.  The one-step safe action set for state $x$ is
\begin{equation}
\mathcal U_{\mathrm{safe}}(x)=[-u_{\max},u_{\max}]\cap[-1-\lambda x,1-\lambda x].
\end{equation}
For $x=r$ and $x=-r$, these sets are individually nonempty but disjoint whenever
\begin{equation}
    1<\lambda r<1+u_{\max}.
\end{equation}
The strict upper bound makes each known state individually controllable.  FDI
creates indistinguishability, so no single history-dependent action certifies
both states.

This example distinguishes fixed filtering from certified learning.  A fixed
severe-attack filter can reject when the common kernel is empty.  A certified
online learner must also prevent updates under a certificate with no feasible
kernel.

\subsection{Certificate-Gated Update}

The resulting update rule first obtains $X_{\mathcal A}(h)$ from a
severe-attack-aware estimator or safety filter.  It then computes
$K_\Gamma(h)$.  An empty kernel rejects the certificate and freezes the
update.  Otherwise, the mechanism accepts, projects, constrains, or rejects the
candidate according to the updated action's kernel membership.

This rule separates fixed-controller safety from the lifecycle condition for
a controller that learns from attacked histories.  The underlying certificate
continues to supply the controller-synthesis and safety obligations.
